\section{Limitations}
\label{sec:limitations}

We are deliberately explicit about what the present results do and do not
establish.

\paragraph{L1 --- Retention is a proxy for end-task accuracy.}
Our metric is gold-evidence retention under a keep budget, not downstream
QA correctness. High retention is \emph{necessary} for a correct answer
that depends on the evidence, but not \emph{sufficient}: the answerer
must also retrieve and use the retained turn. Measuring final QA accuracy
under each memory policy requires an answerer LLM and an LLM judge, which
we leave to an API-gated follow-up. The claim here is about what survives
forgetting, not about answer quality.

\paragraph{L2 --- Three of seven factors are inert in our annotator.}
Value alignment, task utility, and usage history are held at $0$ in the
API-free setting (\cref{tab:factors}): they need a value profile, an LLM
judge, and access logs respectively. The blind headline therefore uses
only four live factors, and task utility --- arguably the most direct
expected-utility signal --- is exactly the one an answerer would unlock.
Our result is thus a \emph{conservative, four-factor, API-free} estimate;
the three inert factors might raise performance or, if noisy, lower it, so
we do not claim the full function is strictly better --- only that a
useful value signal is already recoverable with no API budget.

\paragraph{L3 --- Single benchmark and modality.}
We evaluate on LongMemEval-S ($74$ usable cases of the first $80$),
English multi-session chat. Other agentic workloads --- tool use, coding,
web navigation --- have different gold structure and may favour different
factors; we do not yet test them. The synthetic study controls confounds
but is, by design, not naturalistic.

\paragraph{L4 --- The blind goal anchor is one choice among many.}
We proxy the consolidation-time goal by the session's user-turn centroid.
Other query-agnostic anchors (an explicit task description, a running
objective stack) could carry more signal and would change the blind
goal-relevance term. We claim the \emph{regime distinction} is essential,
not that the session centroid is optimal.

\paragraph{L5 --- Linear value and a simple optimiser.}
\Cref{eq:value} is linear, so it cannot represent factor interactions
(e.g.\ emotion mattering only for self-relevant items). The learner is a
random-restart hill-climb, not CMA-ES; although the blind result is
stable across $20$ resampled splits, a stronger optimiser or a nonlinear
value could shift the fitted weights. The reported standard deviations
reflect split variance over a fixed $74$-case pool, not i.i.d.\ sampling
from the full $500$-question benchmark. The natural next step is to
replace the linear form with a neural $g_\theta$ (\cref{sec:method:value})
--- e.g.\ a small multi-layer perceptron over the factor vector --- and
train it with a differentiable learning-to-rank loss (gold turns ranked
above distractors within a case), since the discrete keep/drop objective
itself is non-differentiable. This would let the value capture the factor
interactions the linear form misses, at the cost of the per-factor
interpretability that is currently a feature; on a $74$-case pilot it
would also overfit, so it should accompany the full $500$-case run.
Because every downstream operation consumes the value only through the
scalar $V(m)$, this swap is contained to the scoring function and leaves
the encoding/forgetting/retrieval machinery unchanged.
